\chapter{Background on Machine Learning and Deep Learning}
\label{ch:ml-methods}

\section{Introduction}

This chapter provides a concise overview of the Machine Learning (ML) and Deep Learning (DL) concepts that are relevant to Fault Detection and Diagnosis (FDD) in photovoltaic systems. It introduces the learning paradigms and model families used in this project, and it establishes the necessary terminology to understand the design choices and experimental results presented in the next chapters.

\section{Machine Learning Methods}
\label{sec:ml_background}

\subsection{Support Vector Machines (SVM)}

Support Vector Machines (SVM) are supervised learning algorithms designed for classification and regression tasks. The main objective of SVM is to determine an optimal separating hyperplane that maximizes the margin between different classes in the feature space. For nonlinear problems, kernel functions are used to project the data into a higher-dimensional space where linear separation becomes possible.

Given a training dataset $\{(x_i,y_i)\}_{i=1}^{N}$, the decision function is expressed as:
\begin{equation}
f(x)=\text{sign}(w^Tx+b),
\end{equation}
where $w$ is the normal vector of the hyperplane and $b$ is the bias term. For nonlinear classification, kernel functions such as the radial basis function (RBF) are commonly employed:
\begin{equation}
K(x_i,x_j)=\exp(-\gamma ||x_i-x_j||^2).
\end{equation}

SVM models are widely used in PV fault diagnosis because of their strong generalization capability and robustness in high-dimensional feature spaces.

\subsection{K-Nearest Neighbors (KNN)}

K-Nearest Neighbors (KNN) is a non-parametric supervised learning algorithm that classifies a sample according to the majority class among its $k$ nearest neighbors in the feature space. Unlike model-based algorithms, KNN does not explicitly learn parameters during training and instead relies on similarity measures between samples.

The Euclidean distance is commonly used:
\begin{equation}
d(x_i,x_j)=\sqrt{\sum_{n=1}^{m}(x_{in}-x_{jn})^2},
\end{equation}
where $m$ is the number of features.

KNN is particularly suitable for PV datasets with moderate dimensionality and nonlinear decision boundaries.

\subsection{Decision Trees}

Decision Trees (DT) are hierarchical supervised learning models that recursively partition the feature space into subsets according to decision rules. Each internal node represents a condition on an input feature, while leaf nodes correspond to output classes.

The selection of the optimal split is generally based on impurity measures such as entropy:
\begin{equation}
H(S)=-\sum_{i=1}^{C}p_i\log_2(p_i),
\end{equation}
where $p_i$ is the probability of class $i$ and $C$ is the number of classes.

Decision Trees are interpretable and computationally efficient, making them suitable for fault classification and rule-based diagnostic systems.

\subsection{Random Forest}

Random Forest (RF) is an ensemble learning technique based on the aggregation of multiple decision trees trained on randomized subsets of data and features. Each tree independently produces a prediction, and the final decision is obtained through majority voting.

If $T_k(x)$ denotes the prediction of the $k$-th tree, the final classifier is:
\begin{equation}
\hat{y}=\text{mode}\{T_1(x),T_2(x),...,T_n(x)\}.
\end{equation}

Random Forest improves robustness and reduces overfitting compared to single decision trees, particularly for noisy PV operational datasets.

\section{Deep Learning Methods}
\label{sec:dl_background}

\subsection{Autoencoders and Variational Autoencoders}

Autoencoders (AE) are unsupervised neural networks designed to learn compressed latent representations of input data. An autoencoder consists of an encoder function that maps the input into a latent space and a decoder function that reconstructs the original input.

The latent representation is obtained as:
\begin{equation}
z=f(Wx+b),
\end{equation}
while the reconstructed output is:
\begin{equation}
\hat{x}=g(W'z+b').
\end{equation}

The network is trained by minimizing the reconstruction error:
\begin{equation}
L=||x-\hat{x}||^2.
\end{equation}

Variational Autoencoders (VAE) extend conventional autoencoders by learning probabilistic latent distributions instead of deterministic representations. Rather than encoding a single latent vector, the encoder estimates the mean $\mu$ and variance $\sigma^2$ of a probability distribution:
\begin{equation}
z \sim \mathcal{N}(\mu,\sigma^2).
\end{equation}

VAEs improve latent space continuity and are suitable for anomaly detection and generative modeling.

\subsection{Long Short-Term Memory (LSTM)}

Long Short-Term Memory (LSTM) networks are recurrent neural networks specifically designed to model temporal dependencies in sequential data. LSTM architectures introduce memory cells and gating mechanisms to overcome the vanishing gradient problem encountered in conventional RNNs.

The forget gate is expressed as:
\begin{equation}
f_t=\sigma(W_f[h_{t-1},x_t]+b_f),
\end{equation}
while the cell state update is:
\begin{equation}
C_t=f_tC_{t-1}+i_t\tilde{C}_t.
\end{equation}

LSTM networks are widely applied in PV fault diagnosis because they effectively capture temporal variations in voltage, current, irradiance, and temperature signals.

\subsection{Transformer Networks}

Transformer models are deep learning architectures based entirely on attention mechanisms rather than recurrent structures. Their main advantage lies in the ability to capture long-range dependencies while enabling parallel computation.

The self-attention mechanism is defined as:
\begin{equation}
\text{Attention}(Q,K,V)=\text{softmax}\left(\frac{QK^T}{\sqrt{d_k}}\right)V,
\end{equation}
where $Q$, $K$, and $V$ denote the query, key, and value matrices.

Transformers are increasingly used in PV fault diagnosis and forecasting due to their strong sequence modeling capability.

\subsection{Convolutional Neural Networks (CNN)}

Convolutional Neural Networks (CNN) are deep architectures specialized for extracting spatial features from structured data such as images or transformed signal matrices. CNNs employ convolutional filters to learn local feature patterns automatically.

The convolution operation is expressed as:
\begin{equation}
S(i,j)=(X*K)(i,j)=\sum_m\sum_n X(m,n)K(i-m,j-n),
\end{equation}
where $X$ is the input matrix and $K$ is the convolution kernel.

CNNs are extensively used in PV systems for thermal image analysis, I--V curve classification, and defect localization because of their strong feature extraction capability.